\documentclass[12pt]{article}
\usepackage[margin=1in]{geometry}
\usepackage{graphicx}
\usepackage{listings}
\usepackage{xcolor}
\usepackage{amsmath}
\usepackage{booktabs}
\usepackage{hyperref}
\usepackage{caption}
\usepackage{float}
\usepackage{amssymb}
\usepackage{textcomp}
\usepackage{titlesec}
\usepackage{parskip}

% Code listing style
\lstdefinestyle{ros}{
    backgroundcolor=\color{gray!10},
    basicstyle=\ttfamily\small,
    keywordstyle=\color{blue},
    commentstyle=\color{green!50!black},
    stringstyle=\color{red!70!black},
    breaklines=true,
    frame=single,
    numbers=left,
    numberstyle=\tiny\color{gray},
    tabsize=4,
    showstringspaces=false
}

\lstdefinestyle{bash}{
    backgroundcolor=\color{black!90},
    basicstyle=\ttfamily\small\color{white},
    breaklines=true,
    frame=single,
    showstringspaces=false
}

\title{
    \vspace{-1cm}
    \large RBE 500 --- Foundations of Robotics \\[0.4em]
    \LARGE \textbf{Project Assignment 3} \\[0.4em]
    \large Worcester Polytechnic Institute
}
\author{Tracy Yang \and Kyle Yee}
\date{\today}

\begin{document}

\maketitle

% ===========================================================================
\section{Velocity-Level Kinematics (Part 1)}
% ===========================================================================

\subsection{Overview}

For part 1, a new node was implemented that serves two functions: 
taking joint velocities and converting them to end effector velocities, 
and taking end effector velocities and converting them to joint velocities.

\subsection{Service Definitions}

Two new services were created in the \texttt{velocity\_kinematics} package.

\begin{lstlisting}[style=ros, caption={JointToEEVelocity.srv}]
float64 theta1_dot
float64 theta2_dot
float64 d3_dot
---
float64 vx
float64 vy
float64 vz
bool success
string message
\end{lstlisting}

\begin{lstlisting}[style=ros, caption={EEToJointVelocity.srv}]
float64 vx
float64 vy
float64 vz
---
float64 theta1_dot
float64 theta2_dot
float64 d3_dot
bool success
string message
\end{lstlisting}

\subsection{Velocity Kinematics Node}

\begin{lstlisting}[style=ros, language=Python,
    caption={velocity\_kinematics\_node.py}]
import rclpy
import numpy as np
from rclpy.node import Node
from sensor_msgs.msg import JointState
from velocity_kinematics.srv import JointToEEVelocity, EEToJointVelocity

A1 = 0.40
A2 = 0.30


def jacobian(theta1, theta2):
    s1, c1 = np.sin(theta1), np.cos(theta1)
    s12, c12 = np.sin(theta1 + theta2), np.cos(theta1 + theta2)
    return np.array([
        [-A1*s1 - A2*s12, -A2*s12, 0],
        [ A1*c1 + A2*c12,  A2*c12, 0],
        [ 0,               0,      1]
    ])


class VelocityKinematicsNode(Node):

    def __init__(self):
        super().__init__('velocity_kinematics_node')
        self.theta1 = 0.0
        self.theta2 = 0.0

        self.create_subscription(JointState, '/joint_states', self.joint_state_cb, 10)
        self.create_service(JointToEEVelocity, '/scara/joint_to_ee_velocity', self.fwd_vel_cb)
        self.create_service(EEToJointVelocity, '/scara/ee_to_joint_velocity', self.inv_vel_cb)

    def joint_state_cb(self, msg):
        if 'joint_1' in msg.name:
            self.theta1 = msg.position[msg.name.index('joint_1')]
        if 'joint_2' in msg.name:
            self.theta2 = msg.position[msg.name.index('joint_2')]

    def fwd_vel_cb(self, request, response):
        qdot = [request.theta1_dot, request.theta2_dot, request.d3_dot]
        vx, vy, vz = jacobian(self.theta1, self.theta2) @ qdot
        response.vx, response.vy, response.vz = vx, vy, vz
        response.success = True
        return response

    def inv_vel_cb(self, request, response):
        xdot = [request.vx, request.vy, request.vz]
        Jv = jacobian(self.theta1, self.theta2)
        q1, q2, q3 = np.linalg.solve(Jv, xdot)
        response.theta1_dot, response.theta2_dot, response.d3_dot = q1, q2, q3
        response.success = True
        return response


def main(args=None):
    rclpy.init(args=args)
    node = VelocityKinematicsNode()
    rclpy.spin(node)
    node.destroy_node()
    rclpy.shutdown()


if __name__ == '__main__':
    main()
\end{lstlisting}

In the velocity kinematics node, there are two main functions: forward velocity callback and inverse velocity callback. 
These functions take theta 1 and theta 2 from a joint state subscriber and calculate the jacobian using a separate function.
With the jacobian, the forward velocity function multiplies it with qdot to get vx, vy, and vz.
The inverse velocity function solves for qdot in the same equation. This equation can fail if the determenant of Jv is 0.

% ===========================================================================
\section{Extending the Position Controller (Part 2)}
% ===========================================================================

\subsection{Overview}

In part 2, the URDF and controller node were updated to extend the position controller from part 2.

\subsection{URDF Changes}

The joint types for joints 1 and 2 were changed back from fixed to revolute. Then, the Gazebo plugin block was updated with joints 1 and 2.

\begin{lstlisting}[style=ros, language=XML,
    caption={joint\_1 and joint\_2 reverted to revolute}]
<joint name="joint_1" type="revolute">
    <parent link="base_link"/><child link="link_1"/>
    <origin xyz="0 0 0.5"/><axis xyz="0 0 1"/>
    <limit effort="1000" lower="-1.57" upper="1.57" velocity="0.5"/>
</joint>

<joint name="joint_2" type="revolute">
    <parent link="arm_1"/><child link="link_2"/>
    <origin xyz="0.4 0 0"/><axis xyz="0 0 1"/>
    <limit effort="1000" lower="-1.57" upper="1.57" velocity="0.5"/>
    <dynamics damping="0.2" friction="0.1"/>
</joint>
\end{lstlisting}

\begin{lstlisting}[style=ros, language=XML,
    caption={Extended Gazebo plugin block}]
<gazebo>
    <plugin filename="gz-sim-joint-state-publisher-system"
            name="gz::sim::systems::JointStatePublisher">
        <topic>/world/empty/model/scara/joint_state</topic>
        <joint_name>joint_1</joint_name>
        <joint_name>joint_2</joint_name>
        <joint_name>joint_3</joint_name>
    </plugin>

    <plugin filename="gz-sim-joint-position-controller-system"
            name="gz::sim::systems::JointPositionController">
        <joint_name>joint_1</joint_name>
        <topic>/scara/joint1_cmd</topic>
        <p_gain>1</p_gain><i_gain>0</i_gain><d_gain>0</d_gain>
    </plugin>
    <plugin filename="gz-sim-joint-position-controller-system"
            name="gz::sim::systems::JointPositionController">
        <joint_name>joint_2</joint_name>
        <topic>/scara/joint2_cmd</topic>
        <p_gain>1</p_gain><i_gain>0</i_gain><d_gain>0</d_gain>
    </plugin>
    <plugin filename="gz-sim-joint-position-controller-system"
            name="gz::sim::systems::JointPositionController">
        <joint_name>joint_3</joint_name>
        <topic>/scara/joint3_cmd</topic>
        <p_gain>1</p_gain><i_gain>0</i_gain><d_gain>0</d_gain>
    </plugin>
</gazebo>
\end{lstlisting}

\subsection{Controller Node Changes}

To reuse the same position controller from part 2, it had to be adjusted to work for multiple joints. 
This was done by creating dictionaries for each of the stored values. Now, each of the variables 
that were originally just for a singular value have values for each of the joints.

\begin{lstlisting}[style=ros, language=Python,
    caption={Per-joint state, publishers, and services}]
JOINT_NAMES = ['joint_1', 'joint_2', 'joint_3']

KP = {'joint_1': 50.0, 'joint_2': 50.0, 'joint_3': 800.0}
KD = {'joint_1': 5.0,  'joint_2': 5.0,  'joint_3': 80.0}
JOINT_MIN = {'joint_1': -1.57, 'joint_2': -1.57, 'joint_3': -0.20}
JOINT_MAX = {'joint_1':  1.57, 'joint_2':  1.57, 'joint_3':  0.20}

self.ref_position = {j: 0.0 for j in JOINT_NAMES}
self.cur_position = {j: 0.0 for j in JOINT_NAMES}
self.prev_error   = {j: 0.0 for j in JOINT_NAMES}

self.cmd_pub = {
    j: self.create_publisher(Float64, f'/scara/{j.replace("_","")}_cmd', 10)
    for j in JOINT_NAMES
}

self.srv = {
    j: self.create_service(
        SetJointPosition,
        f'/scara/set_{j.replace("_","")}_position',
        self.make_set_position_cb(j)
    )
    for j in JOINT_NAMES
}
\end{lstlisting}

\begin{lstlisting}[style=ros, language=Python,
    caption={Callback factory --- binds a joint name to its own service callback}]
def make_set_position_cb(self, joint_name):
    def set_position_cb(request, response):
        target = request.position
        if not (JOINT_MIN[joint_name] <= target <= JOINT_MAX[joint_name]):
            response.success = False
            response.message = f'Requested {target:.4f} outside limits for {joint_name}.'
            return response
        self.ref_position[joint_name] = target
        response.success = True
        response.message = f'{joint_name} reference set to {target:.4f}'
        return response
    return set_position_cb
\end{lstlisting}

\begin{lstlisting}[style=ros, language=Python,
    caption={PD control loop, generalized to all three joints}]
def control_loop(self):
    now = self.get_clock().now()
    if self.prev_time is None:
        self.prev_time = now
        return
    dt = (now - self.prev_time).nanoseconds * 1e-9
    if dt <= 0.0:
        return
    self.prev_time = now

    for j in JOINT_NAMES:
        error = self.ref_position[j] - self.cur_position[j]
        error_dot = (error - self.prev_error[j]) / dt
        self.prev_error[j] = error

        effort = KP[j] * error + KD[j] * error_dot
        effort = max(-EFFORT_LIMIT, min(EFFORT_LIMIT, effort))

        msg = Float64()
        msg.data = effort
        self.cmd_pub[j].publish(msg)
\end{lstlisting}

\end{document}